\section{幂次发散与格点 QCD 中准分布的重正化}

Rossi 与 Resta~\cite{Rossi:2017muf} 对因子化定理 Eq.~(\ref{eq:fact}) 在格点 QCD 中的实际实现
提出了疑问。他们特别考虑其文中式 (19) 的准分布，并试图在强子静止系中计算它的二阶矩。他们指出，
由于迹项与更低量纲算符的混合，被 $1/p_z^2$ 压低的项含有幂次发散的贡献（如其文中式 (35) 所示）。
因此他们断言：``除非对幂次发散项作非微扰减除，否则 $p_z\rightarrow \infty$ 极限不存在。''
然而，仔细的分析表明，在目前的形式体系中这个问题并不存在。

首先，用于抽取物理部分子分布的准分布是 $y=k_z/p_z$ 的函数。形式上可以对它取二阶矩并证明结果
是带双重导数的定域算符。但在实践中，准分布的矩并不存在。正如单圈计算~\cite{Xiong:2013bka}
所示，结果在大 $y$ 处表现为 $1/|y|$，因此除零阶矩之外的所有更高阶矩都不收敛。恰恰是这个大 $y$
发散要求对定域算符作额外重正化；而不取矩时该发散并不出现。

事实上，正是部分子物理的这种非局部准分布表述避免了老式矩计算方法中的幂次发散问题。非局部表述
的紫外（UV）物理简单得多，因而更容易控制发散。

例如，在``重夸克''形式中准分布的重正化是直接的。可以证明，对夸克非单态而言，除威尔逊线自能或
质量修正外不存在其他幂次发散~\cite{Ishikawa:2016znu,Chen:2016fxx,Ji:2017oey,Ishikawa:2017faj,Green:2017xeu}；
其余发散都只是格距的对数发散。因此，准分布可以通过分离出幂次发散的自能因子和对数型重正化因子来
完成重正化，并在格距趋于零时具有光滑极限。注意，若不先分离幂次发散的自能因子，形如
$\alpha_s C_1|z|/a$（$a$ 为格距、$|z|$ 为威尔逊线长度）的单圈质量修正会在双圈于准分布中产生有限
的共线发散项，看起来会破坏准分布的红外因子化~\cite{Li:2016amo}。然而，只要在把格距趋于零之前
先把线性发散从准分布中分离出去，这些项就没有贡献。

用适当方案把准分布重正化之后，就可以利用 Eq.~(\ref{eq:fact}) 的因子化公式把它匹配到
$\overline{\text{MS}}$ 部分子分布；于是高扭度 $1/p_z^2$ 项在 $p_z\rightarrow \infty$ 时光滑消失。
这正是近期文献~\cite{Alexandrou:2017huk,Chen:2017mzz,Stewart:2017tvs} 中 RI/MOM 方案所做的。
